% =============================================================================
% METHOD ONLY - replacement for the "Post-Processing." paragraph in Sec. III-A of Learning_Locomotion_from_Vision__ICRA2027_ (6).pdf
% Two-stage processing as implemented in this repo (scripts/exp1_filter.py: stage A depth outlier rejection, stage B zero-phase 6 Hz Butterworth on the 3-D keypoints) - the processing behind the video row of Fig. 7.
% No evaluation numbers here; the results paragraph is in paper_postprocessing_results.tex.
% Only the "0.5 m jump rule + 3-frame moving average" sentences are replaced; the following "Foot clearance ..." and "Motion Retargeting." paragraphs stay as they are.
% THRESHOLD: option 1 states a fixed 7 cm deviation from the window median; the code currently uses max(3*1.4826*MAD, 5 cm). Both flag the SAME 7 of 1865 valid depth samples on the eight flat clips, i.e. identical filtered trajectories and identical figures - so pick option 1 and change the rule in 3d_recon.py to the fixed threshold, or pick option 2 and keep the code as it is.
% Note for the authors: the expert datasets behind the reported policy results were generated WITHOUT these two stages. Either regenerate them or state that the processing was added after those experiments.
% Units use \, (no siunitx needed); replace by \SI{}{} if the paper loads it.
% =============================================================================

% OPTION 1 - recommended (fixed threshold, no MAD machinery to explain)
\textbf{Post-Processing.} The depth values at the tracked keypoints are noisy, and they become invalid when a keypoint is occluded or the depth measurement returns no reading. We therefore post-process the reconstruction in two stages. First, we reject outliers in the depth measurements: for every keypoint, we compare its depth to the median over a window of $\pm 2$ frames and discard the sample if it deviates by more than $7\,\mathrm{cm}$. Rejected samples, occluded keypoints and missing depth readings are then filled by linear interpolation between the closest valid measurements before and after the gap. Second, we smooth the lifted 3D keypoint trajectories in the global frame using a zero-phase low-pass filter. We use a fourth-order Butterworth filter with a cutoff of $6\,$Hz that we apply forward and backward, so that no delay is introduced and the foot-contact timing is not shifted, in contrast to a causal moving average. The cutoff lies above the stride characteristics of the demonstrations and below the frequency at which the keypoint spectrum flattens into a noise floor, and thus separates the motion from the sensor noise.

% OPTION 2 - same paragraph with the adaptive threshold that the code implements (Hampel filter); use if 3d_recon.py is left unchanged
\textbf{Post-Processing.} The depth values at the tracked keypoints are noisy, and they become invalid when a keypoint is occluded or the depth measurement returns no reading. We therefore post-process the reconstruction in two stages. First, we reject outliers in the depth measurements with a Hampel filter: we discard the depth $d_t$ of a keypoint if it deviates from the median $\tilde{d}$ over a window of $\pm 2$ frames by more than $\max(3\hat{\sigma},\, 5\,\mathrm{cm})$, where $\hat{\sigma} = 1.4826\,\mathrm{med}|d - \tilde{d}|$ is the robust standard deviation of that window. Rejected samples, occluded keypoints and missing depth readings are then filled by linear interpolation between the closest valid measurements before and after the gap. Second, we smooth the lifted 3D keypoint trajectories in the global frame using a zero-phase low-pass filter. We use a fourth-order Butterworth filter with a cutoff of $6\,$Hz that we apply forward and backward, so that no delay is introduced and the foot-contact timing is not shifted, in contrast to a causal moving average. The cutoff lies above the stride characteristics of the demonstrations and below the frequency at which the keypoint spectrum flattens into a noise floor, and thus separates the motion from the sensor noise.

% OPTIONAL add-on clauses (method-level; drop in where useful)
% (a) clip borders - the current code extrapolates there, which produces outliers when a paw is missing at the first or last frame: "; at the clip borders we hold the nearest valid value instead of extrapolating"
% (b) placement of the two stages in the pipeline (Fig. 2): The first stage acts directly on the depth look-up; the second stage acts after the global coordinate transformation, i.e. on metric 3D positions, and before motion retargeting.
% (c) if the illustration of the filter in Fig. 7 should be pointed at from the method rather than from the results: Figure 7 shows the tracked paw height before and after this processing step.
